\section{Candidate scoring}
\label{sec:scoring}

\begin{lede}
Given the moves a robot is \emph{allowed} to make, this is the single number
that ranks them. Every preference in the system --- direction, congestion,
turning cost, aisle rules --- is finally cashed out here, and nowhere else.
\end{lede}

Once PIBT (\Cref{sec:pibt}) has produced a \emph{legal} candidate set, something
must rank it: that ranking is what turns ``collision-free'' into
``collision-free \emph{and} purposeful''. It is entirely separate from
legality. Scoring only reorders candidates that already passed every rejection
rule in \Cref{sec:pibt-rejection}.

\subsection{Proximity mode}
\label{sec:score-mode}

A robot's behaviour should change as it nears its target: far away, only
progress and direction matter; up close, it should not zigzag chasing a
one-step-shorter path. Route distance $r_i = \Dg{w}(\pos{i}{t})$ is bucketed
against $R_{\mathrm{near}} = 2$ and $R_{\mathrm{far}} = 8$:
\begin{equation}
  \mathrm{mode}(r_i) =
  \begin{cases}
    \textsc{transit} & r_i > R_{\mathrm{far}},\\
    \textsc{approach} & R_{\mathrm{near}} < r_i \leq R_{\mathrm{far}},\\
    \textsc{arrival} & r_i \leq R_{\mathrm{near}}.
  \end{cases}
\end{equation}

\subsection{Smooth direction weight \texorpdfstring{$\beta_i$}{beta}}
\label{sec:score-beta}

Rather than a step function on the mode, the weight on ``this move matched my
preferred direction'' decays smoothly as the robot approaches:
\begin{equation}
  \beta_i = \beta_{\max} \cdot
    \min\!\left(1,\ \frac{\max(0,\ r_i - R_{\mathrm{near}})}
      {\max(1,\ R_{\mathrm{far}} - R_{\mathrm{near}})}\right),
  \qquad \beta_{\max} = 3.0.
\end{equation}
At $r_i \geq R_{\mathrm{far}}$ it is at full strength; at
$r_i \leq R_{\mathrm{near}}$ it is zero, leaving the robot free to approach
from any angle. It is $\beta_{\max}$ unconditionally when $r_i = \INF$ (an
unreachable route: keep whatever preference was already computed rather than
collapsing it), and $0$ when \param{direction\_control} is \code{"none"}.

\subsection{Aisle-continuity weight \texorpdfstring{$\gamma_i$}{gamma}}
\label{sec:score-gamma}

A discrete analogue of the same idea, rewarding staying in the current aisle
rather than cutting toward an intersection early:
\begin{equation}
  \gamma_i =
  \begin{cases}
    2.0 & \textsc{transit},\\
    1.25 & \textsc{approach},\\
    0.5 & \textsc{arrival},\\
    0 & \param{direction\_control} = \code{"none"}.
  \end{cases}
\end{equation}

\subsection{Turning cost}
\label{sec:score-turn}

\begin{equation}
  P_{\mathrm{turn}}(\text{prev}, \text{move}) =
  \begin{cases}
    0 & \text{turning cost disabled, or either move is \textsc{stay}},\\
    0 & \text{move} = \text{prev} \quad \text{(straight on)},\\
    \lambda_{\mathrm{reverse}} = 2.0 & \text{move} = \overline{\text{prev}} \quad (180^\circ),\\
    1 & \text{otherwise} \quad (90^\circ).
  \end{cases}
\end{equation}

\subsection{Progress}
\label{sec:score-progress}

The core ``am I getting closer'' signal, with two normalisations
(\param{progress\_normalization}). Writing $d_{\mathrm{cur}} =
\Dg{w}(\pos{i}{t})$, $d_{\mathrm{cand}} = \Dg{w}(v)$ and
$\Delta = d_{\mathrm{cur}} - d_{\mathrm{cand}}$:
\begin{equation}
  \mathrm{progress} =
  \begin{cases}
    \Delta / \max(1, d_{\mathrm{cur}}) & \text{\code{"route"} (the literal specification)},\\
    \Delta \in \{-1, 0, +1\} & \text{\code{"step"} (the default)}.
  \end{cases}
\end{equation}
If the candidate is unreachable, progress is $-1$; if the current position is
unreachable but the candidate is not, progress is $0$ --- neutral rather than
an artificially huge gain.

\begin{caveat}
\code{"step"} is the default rather than the literal formula because at large
$d_{\mathrm{cur}}$ the \code{"route"} normalisation shrinks
$\alpha \cdot \mathrm{progress}$ below $\beta_{\max}$ and
$\lambda_{\mathrm{turn}}$: at $d_{\mathrm{cur}} = 20$,
$\alpha \cdot \mathrm{progress} \approx 10/20 = 0.5 < 3.0 = \beta_{\max}$. The
dominance ordering the design claims (\Cref{sec:score-full}) silently inverts,
and a distant robot prefers keeping its heading to closing distance.
\end{caveat}

\subsection{The full candidate score}
\label{sec:score-full}

\begin{align}
  \score = \ & \alpha \cdot \mathrm{progress} \nonumber\\
    &+ \beta_i \cdot \ind{\text{move} \neq \textsc{stay} \wedge \text{move} = \text{preferred}} \nonumber\\
    &+ \gamma_i \cdot \ind{\text{$v$ stays in the current aisle, robot not at an intersection}} \nonumber\\
    &- \lambda_{\mathrm{turn}} P_{\mathrm{turn}}
      - \mu\, C_i(v)
      - \nu\, P_{\mathrm{wait}}
      - \xi \cdot \ind{\text{$v$ is a bottleneck vertex}} \nonumber\\
    &- \zeta_{\mathrm{counterflow}} \cdot \ind{\text{move opposes the aisle's committed direction}} \nonumber\\
    &- \zeta_{\mathrm{reservation}} \cdot \ind{\text{entering a directional aisle without a ticket}}.
  \label{eq:score}
\end{align}

Defaults: $\alpha = 10.0$, $\lambda_{\mathrm{turn}} = 0.5$, $\mu = 1.0$,
$\nu = 0.2$, $\xi = 1.0$, $\zeta_{\mathrm{counterflow}} =
\zeta_{\mathrm{reservation}} = 8.0$. $P_{\mathrm{wait}} = 1$ unless the
candidate is ``stay in place while already at the waypoint'', which is free;
otherwise waiting is always mildly discouraged, so a robot with any legal
forward move prefers it.

The intended ordering is
\begin{equation}
  \alpha \;>\; \zeta \;>\; \beta_{\max} \;>\; \gamma_{\max} \;>\; \lambda_{\mathrm{turn}}
  \qquad (10 > 8 > 3 > 2 > 0.5):
  \label{eq:ordering}
\end{equation}
a whole step of progress beats everything; driving the wrong way down an aisle
beats every other preference; a single turn is the cheapest term. Two
conditions have to hold for that ordering to be real, and both were violated in
earlier versions of this system.

\paragraph{$\zeta < \alpha$, and $\zeta$ a penalty rather than a rejection.}
The two $\zeta$ terms are where the aisle layer acts. Until this pass they were
hard rejections in \code{feasible\_candidates}, which deletes legal moves from
the candidate set --- and PIBT's progress argument rests on being able to push
any robot into an adjacent cell (\Cref{sec:pibt-recursion}). Pricing counterflow
keeps that freedom: a robot drives the wrong way when that is the only way to
make progress, or when nothing else is left, and pays for it.
\param{hard\_direction\_constraints} restores the old behaviour for comparison,
and \Cref{sec:contribution-soft} is the full argument.

\paragraph{$\mu\, C_i(v) \leq \mu$.} $C_i(v)$ has to be normalised
(\Cref{sec:congestion-mixture}) or the mixture is a robot \emph{count}, and
$\mu C$ then reaches the scale of $\alpha \Delta$ --- measured at mean $3.40$,
$p_{90}$ $5.75$, max $9.60$ on \variant{warehouse\_medium} with 40 robots. That
makes congestion the second-largest term in the score rather than a modulator
of the smaller ones.

\code{CandidateScorer.sort\_key(robot, v)} is $(-\score, v)$: descending score,
ties broken by vertex coordinates for determinism.

% worked-example: score
\begin{workedexample}
\textbf{Worked example.} \code{warehouse\_medium}, variant \code{full\_lda\_pibt}, seed 7, 24 robots, timestep 18. Robot 3 is at (10,20) heading for waypoint (0,12), 18 steps away, so it is in \code{TRANSIT} mode with \code{β\_i} = 3 and \code{γ\_i} = 2 (\Cref{sec:score-beta}, \Cref{sec:score-gamma}). Its previous move was \code{south} and its preferred direction is \code{north}. Every column below is computed by the same functions \code{CandidateScorer.score} itself calls, and they sum to the printed \code{S(v)}:

{\fontsize{5.6}{6.7}\selectfont
\begin{verbatim}
cell     move   α·prog  β·dir  γ·aisle  −λ·turn  −μ·cong  −ν/−ξ  −ζ   S(v)    rejected by      priced
-------  -----  ------  -----  -------  -------  -------  -----  ---  ------  ---------------  -------------------------------
(9,20)   north  +10     +3     +0       −1       −0.11    +0     +0   11.89   vertex-conflict  -
(10,19)  west   +10     +0     +0       −0.5     −0.1     +0     +0   9.4     -                -
(10,20)  stay   +0      +0     +0       +0       −0.16    −0.2   +0   -0.36   -                -
(10,21)  east   −10     +0     +0       −0.5     −0.16    +0     +0   -10.66  -                -
(11,20)  south  −10     +0     +0       +0       −0.28    +0     −16  -26.28  -                aisle-direction, no-reservation
\end{verbatim}
}

Read the last two columns first, because they are different kinds of thing. \textbf{Rejected by} is legality: those moves are gone before scoring, and no score could rescue them. \textbf{Priced} is preference: the move is still on the table, it just costs.

(9,20) scores highest at 11.89 and is unavailable regardless: \code{vertex-conflict} removed it in \Cref{sec:pibt-rejection}. So the move actually chosen is (10,19) at S = 9.4 — the best of what is \emph{legal}.

The shape of the winning number is the whole argument: progress contributes +10, while every other term together moves it by −0.6. The preference terms break ties between moves that agree on progress; they never outvote progress itself. That is what \code{α > ζ > β > γ > λ} buys.

(11,20) is the instructive row. It pays 16 for \code{aisle-direction, no-reservation} and finishes 35.68 behind — last by a wide margin, and still in the candidate set. Set \code{hard\_direction\_constraints} and that same move is deleted instead; a move PIBT cannot see is a move priority inheritance cannot use to unblock a chain (\Cref{sec:pibt-recursion}). The gap between those two treatments is what the README's ablation measures as 1.9× to 3.1× throughput.
\end{workedexample}
% /worked-example

\subsection{Preferred route direction and hysteresis}
\label{sec:score-hysteresis}

\begin{algorithm}[H]
\caption{\code{compute\_route\_direction(robot)}}
\begin{algorithmic}[1]
\If{$|\mathrm{route}| \geq 2$}
  \Return the compass direction from $\mathrm{route}[0]$ to $\mathrm{route}[1]$
\Else
  \ \Return the neighbour $n$ maximising $\Dg{w}(\pos{i}{t}) - \Dg{w}(n)$
  \Comment{steepest single-step descent}
\EndIf
\end{algorithmic}
\end{algorithm}

\code{apply\_direction\_hysteresis} keeps the \emph{previous} preferred
direction unless one of: hysteresis is off; the previous preference was
\textsc{stay}; the proposed direction already matches; the robot has been
blocked for at least $t_{\mathrm{blocked}}$ steps; the robot is not in
\textsc{transit} mode; or the robot is at an intersection. In other words,
\textbf{new direction commitments only happen at decision points}, which stops
a robot flip-flopping mid-corridor as the greedy choice above wobbles.
